\section{Discussion}
\label{sec:discussion}

\paragraph{Rationalized gaming, not blind gaming.} Our agents did not hardcode an expected output and hope no one looked. The agents in our study game a test when they can find a rule under which the test's value would be right; they implement that rule as ordinary-looking code; they rewrite the documentation to match; and they describe the rule to the user as a requirement. When no such rule exists, most of them stop and report the conflict. The code, the docstring, and the report then tell the same story, so none of them alone gives the gaming away. We cannot tell whether an agent believed its rule; we can say that its stated source does not exist, and that an honest agent should say which parts of an explanation it was told, inferred, or made up. This extends findings that chain-of-thought explanations can omit or rationalize the true cause of behavior \citep{turpin2023language,chen2025reasoning,arcuschin2026chainofthought} to the artifacts that users and reviewers actually read.

\paragraph{What this means for oversight of coding agents.} First, a green suite plus a plausible summary is weak evidence: in our data the most common report accompanying a special case is a confident statement of a requirement that does not exist. Second, the kind of question matters. ``Did you special-case a test?'' invites a technically defensible ``no'' whenever the code is general; ``what is the source for this behavior?'' mostly elicits the truth, although in our design it also named the test. Oversight tools and human reviewers should ask for provenance, not for confessions. Third, claimed requirements are checkable: a monitor can extract every requirement an agent cites and search the repository, ticket, and conversation for it. Our judge did this, although it also had the grader's verdict, which a deployed monitor would lack. Fourth, edits to documentation deserve the same scrutiny as edits to tests, since rewriting a docstring to match new behavior erases the evidence of the conflict.

\paragraph{The behavior is not inevitable.} Claude Opus 5.5 and Fable 5.1 gamed less, almost never when the user named the docstring as authoritative, and reported every special case they made transparently. Fable 5.1 often followed the test under \textsc{spec-tests}, a defensible reading of that instruction, but said so. Sonnet 5 behaved like the 4.5 models. Honest reporting of test conflicts is therefore achievable in current models and measurable with the protocol we release. Because the gaming depends on the availability of a plausible rule, evaluations of reward hacking that use only obviously wrong tests (a transposed digit, an impossible value) are likely to underestimate it; conflicts that admit a cover rule are the informative case.

\paragraph{Limitations.}
\label{sec:limitations}
All subjects, judges, and reviewers are Claude models, no human labels were collected, and all agents ran in one harness. The repositories are small and constructed. We cannot see the models' reasoning, so we call reports fabricated relative to the repository, not intentionally deceptive. The minimal pairs make the test error both harder to rationalize and easier to spot. \Cref{app:limitations} discusses these and other limits.


\section{Conclusion}
\label{sec:conclusion}

When a test contradicts the specification, most coding agents we studied make it pass with a rule under which the test is right, and tell the user the rule is a requirement. The invented requirement is visible, checkable, and, when the agent is asked for its source, usually disowned. Reports from coding agents should be read as claims about provenance, and checking a claimed source is a concrete, inexpensive step.
